%%%%%%%%%%%%%%%%%%%%%%%%%%%%%%%%%%%%%%%%%%%%%%%%%%%%%%%%%%%%%%%%%%%%%%%%%%%
% Lehreinheit 13: Wissenschaftsethik und Grenzarbeit
% Quelle: WuZ-Buch, Kap. 15 (LE 13), Anhang E Exposure und Grenzarbeit
%         (Kommentar M. Pfeil, Nachtrag Dufourmantelle), Kap. 5c Offenheit
%         und Geheimnis (Merton, Vertrauensspiel), Kap. 5b Vogelmann,
%         Anhang C Magiculum, Kap. 15b Entlarvung (Forschungsbetrug)
% Staunmoment: Boundary Work: Was verrät man?  |  Kernbegriff: Forschungsethik
%%%%%%%%%%%%%%%%%%%%%%%%%%%%%%%%%%%%%%%%%%%%%%%%%%%%%%%%%%%%%%%%%%%%%%%%%%%
\nsl{Die letzte Einheit stellt eine Frage, die das ganze Semester begleitet
hat, ohne dass sie ausgesprochen wurde: Was muss man offenlegen, und was
darf man für sich behalten? Der Zauberkünstler hütet seine Methode, die
Wissenschaft muss ihre offenlegen. Die Einheit beginnt mit einem Trick, der
gezeigt und anschließend offen verhandelt wird, aber nicht erklärt. Aus
dieser Grenzarbeit entwickelt sie die Grundsätze der Forschungsethik:
informierte Einwilligung, Anonymisierung und Datenschutz, wissenschaftliche
Integrität mit den drei Kerndelikten Erfindung, Fälschung und Plagiat, Open
Science, Interessenkonflikte, Ethikkommissionen und die Grenzen dessen, was
Wissenschaft tun soll. In der Ethik-Werkstatt prüfen die Studierenden ihr
eigenes Semesterprojekt. Den Abschluss bildet der Rückblick auf den Bogen
vom Staunen in LE~1 zur Verantwortung in LE~13.
\newpage}
\os{\begin{description}
\item[{\rot\bf Heute}]~\\[-\lblskip]
\bi
  \ite ein Trick -- und die Frage, ob ich ihn verrate
  \ite Grenzarbeit: offenlegen oder schützen?
  \ite Einwilligung, Anonymisierung, DSGVO
  \ite Integrität: Erfindung, Fälschung, Plagiat
  \ite Open Science, Interessenkonflikte, Ethikkommissionen
  \ite Rückblick: vom Staunen zur Verantwortung
\ei
\end{description}
\newpage}

\begin{vorbereitung}[Boundary Work: Was verrät man?]
\textbf{Eine Woche vorher}
\begin{checkliste}
  \item \textbf{Trick wählen.} Naheliegend ist die Reprise des
    Switchboards aus LE~5: Die Gruppe kennt den Effekt, hat in LE~5
    Hypothesen gebildet und ist mit hoher Wahrscheinlichkeit noch neugierig.
    Ebenso geeignet ist jeder andere Effekt des Semesters, dessen Methode
    nicht im Buch offengelegt ist (z.\,B.\ der Seileffekt aus LE~11).
    Methode und Einstudierung nach der Anleitung des Herstellers bzw.\ der
    Standardliteratur; hier bewusst nicht beschrieben.
  \item Den Effekt erneut proben; nach einigen Wochen Pause sitzen
    Handgriffe oft schlechter als gedacht.
  \item Eigene Position zur Exposure-Frage klären und eine Formulierung
    vorbereiten, die die \emph{Richtung} der Methode andeutet, ohne sie zu
    verraten (siehe Abschnitt~\ref{sec:le13-grenzarbeit}).
  \item Klären, welche Stelle an der eigenen Hochschule für
    Ethikfragen, Datenschutz und gute wissenschaftliche Praxis zuständig
    ist (Ethikkommission, Datenschutzbeauftragte, Ombudsperson), und die
    Kontaktdaten auf eine Folie setzen.
  \item Den Studierenden mitteilen: Stand des Semesterprojekts mitbringen,
    insbesondere Art der erhobenen oder geplanten Daten und Teilnehmende.
\end{checkliste}

\textbf{Handout und Übungsblatt}
\begin{checkliste}
  \item Handout \enquote{Wissenschaftsethik und Grenzarbeit} sowie
    Übungsblatt LE~13 Teil~A und Teil~B je einmal pro Person kopieren.
    Teil~A (Checkliste, Einwilligung, Fallvignetten) liegt zu Beginn aus;
    Teil~B (Intransparenz und Täuschung, Abschluss-Tagebucheintrag) wird
    zur Verantwortung nach außen ausgeteilt.
  \item Der Einwilligungs-Mustertext (Aufgabe~2) ist ein Gerüst. Wenn die
    Hochschule eigene Vorlagen hat, diese zusätzlich ausgeben.
  \item Für die Fallvignetten (Aufgabe~3) die neun Fälle zusätzlich auf
    Karten drucken (eine Karte pro Fall), wenn sie in Gruppen sortiert
    werden sollen.
  \item Die Tagebücher des Staunens werden heute abgeschlossen; Abgabe-
    modus für die Portfolioprüfung (20\,\%) ankündigen.
\end{checkliste}

\textbf{Am Tag selbst}
\begin{checkliste}
  \item Trick aufgebaut, aber abgedeckt; die Einheit beginnt mit der
    Vorführung.
  \item Tafel teilen: \enquote{Was verlöre man durch Offenlegung?} und
    \enquote{Was gewönne man?}
  \item Rückblick-Folie (13 Kernbegriffe) für den Schluss bereit; ggf.\
    die Requisiten des Semesters (Pendel, Seil, Karten) sichtbar
    auslegen.
\end{checkliste}

\textbf{Risiken und Plan B.} Die Gruppe wird darauf drängen, die Methode zu
erfahren; manche werden sie bei einer Reprise zu erraten glauben. Beides
ist Material für die Diskussion: Ein geäußerter Lösungsvorschlag wird
weder bestätigt noch verneint, sondern als Hypothese behandelt
(\enquote{Was würde die Beobachtung sonst noch erklären?}). Wer nach dem
Kurs privat nachfragt, bekommt einen Hinweis auf die einschlägige
Fachliteratur, nicht die Methode.
\end{vorbereitung}

\begin{ablauf}
\begin{description}
\item[0--10 · Staunmoment/Diskurs.] Ein Trick wird gezeigt und offen
  verhandelt: Was verlöre man durch vollständige Offenlegung, was gewönne
  man? (Methode vs.\ Erleben.)
\item[10--40 · Theorie-Input.] Grenzarbeit; Forschungsethik (Einwilligung,
  Anonymisierung, DSGVO); Integrität (FFP, DFG-Kodex); Open Science;
  Interessenkonflikte; Ethikkommissionen; Grenzen der Wissenschaft.
\item[40--65 · Ethik-Werkstatt.] Anwendung der DFG-Leitlinien auf das
  eigene Semesterprojekt mit der Checkliste (Aufgabe~1); Skizze eines
  Ethikplans; Fallvignetten (Aufgabe~3); alles auf Übungsblatt Teil~A.
\item[65--80 · Verantwortung nach außen.] Übungsblatt Teil~B austeilen.
  Kurzbrücke zur Desinformationsabwehr: Wann wird Intransparenz zur
  Täuschung? (Aufgabe~4)
\item[80--90 · Semesterabschluss.] Abschluss-Tagebucheintrag
  (Aufgabe~5); Rückblick
  auf den Bogen von LE~1 (Staunen) zu LE~13 (Verantwortung).
\end{description}
\end{ablauf}

%==============================================================================
\section{Staunmoment: Was verrät man?}
%==============================================================================
\begin{staunbox}{Boundary Work}
Sie sehen einen Trick, den manche von Ihnen schon kennen.

\medskip
\emph{Danach entscheiden wir gemeinsam: Soll ich Ihnen die Methode
verraten? Sammeln Sie Gründe dafür und dagegen.}
\end{staunbox}
\os{\newpage}

\begin{description}
\item[{\rot\bf Offenlegen?}]~\\[-\lblskip]
\bi
  \ite Was \emph{verlören} Sie, wenn Sie die Methode kennen?
  \ite Was \emph{gewönnen} Sie?
  \ite Wäre es danach derselbe Trick?
  \ite Abstimmung per Handzeichen: verraten -- nicht verraten
\ei
\os{\newpage}

\item[{\rot\bf Methode und Erleben}]~\\[-\lblskip]
\begin{center}\nsl{\small}
\begin{tabular}{p{0.42\textwidth}p{0.42\textwidth}}
\toprule
\textbf{durch Offenlegung verloren} & \textbf{durch Offenlegung gewonnen} \\
\midrule
das Staunen, die offene Frage & das Wissen, wie es geht \\
der Anlass, weiterzudenken & Sicherheit: kein Übernatürliches \\
der Trick für alle künftigen Vorführungen & ein Puzzle, das gelöst ist \\
\bottomrule
\end{tabular}
\end{center}
\bi
  \ite Die Methode ist kein großes Geheimnis -- das Erleben ist es
\ei
\os{\newpage}
\end{description}

\begin{lehrhinweis}
Die Abstimmung fällt meist für das Verraten aus. Dann nicht verraten, aber
begründen: Die Gruppe hat soeben selbst gesammelt, was verloren ginge, und
die rechte Spalte enthält kaum etwas, das nicht auch die Andeutung leisten
würde (\enquote{Es gibt eine technische Lösung, die im Gerät steckt, und
sie hat nichts mit mir zu tun.}). Markus Pfeil beschreibt im Buch genau
diese Erfahrung mit dem Switchboard: Ohne Hinweis schreiben Studierende den
Effekt dem Vorführenden zu, mit dem Hinweis auf eine technische Lösung
begreifen sie, wie wenig sie von einem scheinbar einfachen System
verstehen. Die Tabelle ist eine erwartbare Sammlung; eigene Punkte der
Gruppe haben Vorrang.
\end{lehrhinweis}

%==============================================================================
\section{Grenzarbeit: offenlegen oder schützen}
\label{sec:le13-grenzarbeit}
%==============================================================================
\nsl{Der Soziologe Thomas Gieryn hat den Begriff \emph{boundary-work}
geprägt: Die Grenze zwischen Wissenschaft und Nicht-Wissenschaft ist keine
feste Eigenschaft des Wissens, sondern wird gezogen, verteidigt und
verschoben \parencite{gieryn1983boundary}. Brian Rappert und Gustav Kuhn
haben den Begriff auf die Enthüllung von Zaubertricks übertragen; Markus
Pfeil hat ihren Ansatz im Anhang des Buchs aus Sicht der Lehre kommentiert.
Seine Position: Methoden werden nicht verraten, weil sie das Geheimnis
tragen; die \emph{Richtung}, in der eine Methode liegt, darf und soll man
andeuten. Mit Anne Dufourmantelle lässt sich das schärfer fassen: Eine
Enthüllung ist keine Subtraktion von Unwissen, sondern verwandelt die
Erfahrung, und das Geheimnis ist eine Handlung, kein Gegenstand
\parencite{dufourmantelle2021geheimnis}.

Die Wissenschaft steht auf der anderen Seite dieser Grenze. Robert K.\
Merton hat ihre Normen 1942 als \emph{CUDOS} zusammengefasst: Communism
(gemeinsames Eigentum an Erkenntnis), Universalism, Disinterestedness,
Organized Skepticism \parencite{merton1973normative}. Wer sein Verfahren
verbirgt, entzieht sich der Prüfung durch andere. Das Buch spricht von
einer \emph{Inversion}: Dieselbe Handlung, die Offenlegung, vollendet die
Wissenschaft und vernichtet die Zauberkunst. Auch im Buch selbst wird die
Grenze unterschiedlich gezogen. Im Kapitel über Offenheit und Geheimnis
wird das Geheimnis im Unterricht \enquote{verzögert, aber stets
eingelöst}, im Kommentar zur Exposure-Frage bleibt die Methode
geheim und nur das Prinzip wird erklärt. Diese Spannung ist selbst ein
Stück Grenzarbeit.}
\begin{description}
\item[{\rot\bf Grenzarbeit}]~\\[-\lblskip]
\bi
  \ite Grenzen zwischen Wissen und Nichtwissen werden gezogen, verteidigt,
    verschoben \parencite{gieryn1983boundary}
  \ite die Wissenschaft selbst entstand durch Grenzarbeit, etwa die
    Trennung der Chemie von der Alchemie \parencite{newmanprincipe2002}
  \ite Position im Buch: Methode nicht verraten, Richtung andeuten
  \ite Enthüllung verwandelt die Erfahrung
    \parencite{dufourmantelle2021geheimnis}
\ei
\os{\newpage}

\item[{\rot\bf Eine Inversion}]~\\[-\lblskip]
\begin{center}
\resizebox{0.9\textwidth}{!}{%
\begin{tikzpicture}[font=\sffamily,>=Stealth,
  b/.style={draw=rwuviolet,thick,fill=rwuviolet!10,rounded corners=3pt,
    minimum height=16mm,align=center,text width=40mm},
  c/.style={b,draw=rwucyan,fill=rwucyan40!40}]
  \node[c] (w) at (0,0) {\textbf{Wissenschaft}\\Offenlegung\\= Vollendung};
  \node[b] (z) at (8,0) {\textbf{Zauberkunst}\\Offenlegung\\= Zerstörung};
  \node[draw=rwugray,dashed,rounded corners=3pt,align=center,text width=56mm,minimum height=10mm] (m) at (4,-3)
    {dieselbe Handlung: das Verfahren zeigen};
  \draw[->,thick] (m) -- (w);
  \draw[->,thick] (m) -- (z);
\end{tikzpicture}}
\end{center}
\bi
  \ite CUDOS: Communism, Universalism, Disinterestedness, Organized
    Skepticism \parencite{merton1973normative}
\ei
\os{\newpage}

\item[{\rot\bf Was in der Forschung offenliegen muss -- und was nicht}]~\\[-\lblskip]
\begin{center}\nsl{\small}
\begin{tabular}{p{0.42\textwidth}p{0.46\textwidth}}
\toprule
\textbf{muss offenliegen} & \textbf{darf/muss geschützt bleiben} \\
\midrule
Methode, Material, Auswertung & Identität der Teilnehmenden \\
Daten (soweit möglich) & sensible Rohdaten (z.\,B.\ Gesundheit) \\
Finanzierung, Interessenkonflikte & Begutachtung (Vertraulichkeit) \\
Abweichungen vom Plan & Rechte Dritter (Lizenzen, Patente) \\
\bottomrule
\end{tabular}
\end{center}
\bi
  \ite Leitlinie: so offen wie möglich, so geschützt wie nötig
\ei
\os{\newpage}
\end{description}

%==============================================================================
\section{Forschung mit Menschen: Einwilligung}
%==============================================================================
\nsl{Wer Menschen befragt, beobachtet oder testet, braucht ihre
\emph{informierte Einwilligung}. Informiert heißt: Die Teilnehmenden wissen
vorher, worum es geht, was sie tun sollen, wie lange es dauert, welche
Risiken bestehen, was mit ihren Daten geschieht und an wen sie sich wenden
können. Einwilligung heißt: Die Teilnahme ist freiwillig und kann jederzeit
ohne Nachteil beendet werden; nach der Datenschutz-Grundverordnung ist eine
Einwilligung jederzeit widerrufbar \parencite{dsgvo2016}. Die Grundsätze
gehen auf die Medizinethik nach 1945 zurück (Nürnberger Kodex, Deklaration
von Helsinki) und gelten heute für alle Disziplinen, die mit Menschen
forschen. Ein Sonderfall betrifft genau dieses Semester: Studien mit
Zaubertricks täuschen ihre Teilnehmenden oft absichtlich. Eine solche
Täuschung ist nur vertretbar, wenn die Frage anders nicht zu beantworten
ist, kein ernsthafter Schaden droht und die Teilnehmenden danach aufgeklärt
werden (\emph{Debriefing}).}
\begin{description}
\item[{\rot\bf Informierte Einwilligung}]~\\
\defn{einwilligung}{\textbf{Informierte Einwilligung} ist die freiwillige,
widerrufbare Zustimmung einer Person zur Teilnahme an einer Studie,
nachdem sie über Zweck, Ablauf, Dauer, Risiken und den Umgang mit ihren
Daten verständlich aufgeklärt wurde.}
\os{\newpage}

\item[{\rot\bf Was in die Information gehört}]~\\[-\lblskip]
\bi
  \ite Wer forscht, wozu, und was soll ich tun -- wie lange?
  \ite Risiken und Belastungen (auch: \enquote{keine bekannten})
  \ite freiwillig, jederzeit abbrechbar, ohne Nachteil
  \ite welche Daten, wie gespeichert, wie lange, wer sieht sie
  \ite Kontakt für Fragen und Widerruf
\ei
\os{\newpage}

\item[{\rot\bf Täuschung in Studien}]~\\[-\lblskip]
\bi
  \ite Studien mit Zaubertricks täuschen absichtlich -- das gehört zur
    Methode
  \ite vertretbar nur, wenn die Frage anders nicht zu beantworten ist und
    kein ernsthafter Schaden droht
  \ite danach \textbf{Debriefing}: aufklären, was untersucht wurde
    (nicht unbedingt, wie der Trick geht)
  \ite Teilnehmende dürfen danach ihre Daten zurückziehen
\ei
\os{\newpage}
\end{description}

%==============================================================================
\section{Anonymisierung, Pseudonymisierung, DSGVO}
%==============================================================================
\nsl{Personenbezogene Daten sind alle Informationen, die sich auf eine
identifizierte oder identifizierbare Person beziehen. Die
Datenschutz-Grundverordnung unterscheidet zwei Formen, den Personenbezug zu
lösen. \emph{Pseudonymisierte} Daten können nur mit einer getrennt
aufbewahrten Zusatzinformation (etwa einer Codeliste) einer Person
zugeordnet werden (Art.~4 Nr.~5); sie bleiben personenbezogen und fallen
unter die Verordnung. \emph{Anonyme} Daten lassen sich auch mit
vertretbarem Aufwand keiner Person mehr zuordnen; für sie gilt die
DSGVO nicht (Erwägungsgrund~26) \parencite{dsgvo2016}. Für ein
Semesterprojekt mit Fragebogen ist das wichtigste Prinzip die
\emph{Datenminimierung}: nur erheben, was für die Frage nötig ist. Wer keine
Namen braucht, erhebt keine. Gesundheitsdaten und andere besondere
Kategorien (Art.~9) sind besonders geschützt und gehören in einem
studentischen Projekt nur nach Rücksprache in die Erhebung.}
\begin{description}
\item[{\rot\bf Anonym oder pseudonym?}]~\\
\defn{pseudonymisierung}{Daten sind \textbf{pseudonymisiert}, wenn sie nur
mit einer getrennt aufbewahrten Zusatzinformation einer Person zugeordnet
werden können; sie bleiben personenbezogen. Daten sind \textbf{anonym},
wenn eine Zuordnung auch mit vertretbarem Aufwand nicht mehr möglich ist
\parencite{dsgvo2016}.}
\bi
  \ite Vorsicht: Alter + Studiengang + Semester kann in einem kleinen Kurs
    eindeutig sein
\ei
\os{\newpage}

\item[{\rot\bf Grundsätze der DSGVO (Art.~5)}]~\\[-\lblskip]
\bi
  \ite \textbf{Rechtmäßigkeit, Transparenz}: Rechtsgrundlage, z.\,B.\
    Einwilligung; offen sagen, was geschieht
  \ite \textbf{Zweckbindung}: nur für den genannten Zweck
  \ite \textbf{Datenminimierung}: nur, was nötig ist
  \ite \textbf{Richtigkeit, Speicherbegrenzung}: korrekt halten, löschen,
    wenn nicht mehr nötig
  \ite \textbf{Integrität, Vertraulichkeit}: sicher speichern
  \ite \textbf{Rechenschaft}: nachweisen können, dass man es so gemacht hat
\ei
\os{\newpage}

\item[{\rot\bf Für Ihr Semesterprojekt}]~\\[-\lblskip]
\bi
  \ite keine Namen, keine E-Mail-Adressen im Datensatz
  \ite falls nötig: Code statt Name, Codeliste getrennt und gesichert
  \ite Speicherung auf Hochschulsystemen, nicht im privaten Cloud-Konto
  \ite Altersgruppen statt Geburtsdatum; kleine Zellen zusammenfassen
  \ite Löschtermin festlegen und einhalten
\ei
\os{\newpage}
\end{description}

%==============================================================================
\section{Wissenschaftliche Integrität}
%==============================================================================
\nsl{Als Kerntatbestände wissenschaftlichen Fehlverhaltens gelten
\emph{Fabrication} (Erfinden von Daten oder Ergebnissen),
\emph{Falsification} (Verfälschen von Daten, Material oder Verfahren) und
\emph{Plagiarism} (Übernahme fremder Leistung ohne Kennzeichnung, LE~11).
Das Buch zeigt alle drei an realen Fällen. Diederik Stapel erfand über
Jahre Datensätze; die Untersuchungskommission stellte fest, dass mindestens
55 Publikationen und 10 Dissertationen auf erfundenen Daten beruhten
\parencite{levelt2012flawed}. Reiner Protsch datierte Fossilien falsch und
gab fremde Arbeiten als eigene aus. Elisabeth Bik sichtete die Abbildungen
aus 20.621 biomedizinischen Aufsätzen und fand in 3,8\,\% problematische
Bildduplikate, bei mindestens der Hälfte davon mit Hinweisen auf bewusste
Manipulation \parencite{bik2016prevalence}. Zwischen redlicher Arbeit und
FFP liegt eine Grauzone fragwürdiger Praktiken, die LE~10 behandelt hat:
nachträgliches Auswählen der \enquote{passenden} Auswertung, Hypothesen
nach Sicht der Daten, selektives Berichten
\parencite{simmons2011falsepositive}. In Deutschland beschreibt der Kodex der
Deutschen Forschungsgemeinschaft in 19 Leitlinien, was gute
wissenschaftliche Praxis ist; Hochschulen müssen ihn umsetzen, um
DFG-Mittel zu erhalten \parencite{dfg2019kodex}.}
\begin{description}
\item[{\rot\bf FFP: drei Kerndelikte}]~\\[-\lblskip]
\begin{center}\nsl{\small}
\begin{tabular}{p{0.24\textwidth}p{0.3\textwidth}p{0.32\textwidth}}
\toprule
\textbf{Delikt} & \textbf{bedeutet} & \textbf{Fall im Buch} \\
\midrule
Fabrication & Daten erfinden & Stapel \parencite{levelt2012flawed} \\
Falsification & Daten, Bilder, Verfahren verfälschen & Protsch; Bildduplikate \parencite{bik2016prevalence} \\
Plagiarism & fremde Leistung ohne Kennzeichnung & Protsch; LE~11 \\
\bottomrule
\end{tabular}
\end{center}
\bi
  \ite dazwischen: fragwürdige Praktiken wie p-Hacking und HARKing (LE~10)
\ei
\os{\newpage}

\item[{\rot\bf Wer hat es aufgedeckt?}]~\\[-\lblskip]
\bi
  \ite Stapel: drei junge Mitarbeitende, denen die Daten zu perfekt waren
  \ite Protsch: eine unabhängige Neudatierung
  \ite Bik: Handarbeit an über 20.000 Aufsätzen
    \parencite{bik2016prevalence}
  \ite das stärkste Argument für die Wissenschaft: Sie deckt ihre Fehler
    selbst auf
\ei
\os{\newpage}

\item[{\rot\bf Der DFG-Kodex (2019)}]~\\[-\lblskip]
\bi
  \ite 19 Leitlinien guter wissenschaftlicher Praxis
    \nsl{\parencite{dfg2019kodex}}
  \ite u.\,a.: lege artis arbeiten, dokumentieren, Daten archivieren,
    Ergebnisse zugänglich machen
  \ite Autorschaft nur für einen genuinen Beitrag
  \ite Vertraulichkeit bei Begutachtungen
  \ite Ombudspersonen: Anlaufstelle bei Verdacht auf Fehlverhalten
\ei
\os{\newpage}
\end{description}

%==============================================================================
\section{Open Science}
%==============================================================================
\nsl{Die Antwort der Wissenschaft auf Betrug und fragwürdige Praktiken ist
weniger Moral als Offenlegung. Das Buch beschreibt Forschung als
Vertrauensspiel: Niemand kann alle Experimente selbst wiederholen, auf die
er sich stützt, also braucht das System Regeln, die zur Offenlegung zwingen.
Die Open-Science-Bewegung bündelt solche Regeln. \emph{Open Access} macht
Publikationen frei zugänglich. \emph{Open Data} macht Daten nachnutzbar,
und zwar nach den FAIR-Prinzipien auffindbar, zugänglich, interoperabel und
wiederverwendbar; FAIR heißt nicht, dass alle Daten offen sein müssen,
sondern dass ihre Bedingungen klar beschrieben sind
\parencite{wilkinson2016fair}. Die \emph{Präregistrierung} legt Hypothesen
und Auswertungsplan vor der Datenerhebung fest und macht so sichtbar, was
geprüft und was nachträglich gefunden wurde
\parencite{nosek2018preregistration}. Wie nötig das ist, zeigte das
Replikationsprojekt der Open Science Collaboration, das nach der
Darstellung im Buch nur etwa ein Drittel bis die Hälfte der untersuchten
psychologischen Befunde bestätigen konnte
\parencite{opensciencecollaboration2015}.}
\begin{description}
\item[{\rot\bf Vier Formen der Offenheit}]~\\[-\lblskip]
\bi
  \ite \textbf{Open Access}: Publikationen frei lesbar
  \ite \textbf{Open Data}: Daten FAIR -- auffindbar, zugänglich,
    interoperabel, wiederverwendbar \parencite{wilkinson2016fair}
  \ite \textbf{Open Material/Code}: Fragebogen, Stimuli, Auswertungsskript
  \ite \textbf{Präregistrierung}: Plan vor den Daten
    \parencite{nosek2018preregistration}
\ei
\os{\newpage}

\item[{\rot\bf Warum Offenheit?}]~\\[-\lblskip]
\bi
  \ite Replikationsprojekt: nur etwa ein Drittel bis die Hälfte der
    geprüften Befunde bestätigt \parencite{opensciencecollaboration2015}
  \ite Präregistrierung ist eine \enquote{glaubwürdige Selbstbindung} --
    sie macht Schönen nachträglich sichtbar
  \ite Offenheit und Datenschutz widersprechen sich nicht: anonymisieren,
    dann teilen
\ei
\os{\newpage}
\end{description}

%==============================================================================
\section{Interessenkonflikte und Ethikkommissionen}
%==============================================================================
\nsl{Mertons Norm der \emph{Disinterestedness} verlangt nicht, dass
Forschende keine Interessen haben, sondern dass Interessen das Ergebnis
nicht steuern. Weil sich das von außen nicht prüfen lässt, gilt die Regel,
Interessenkonflikte offenzulegen: Finanzierung, Beraterverträge,
Beteiligungen, auch persönliche Nähe zum Gegenstand. Der Fall Theranos, den
das Buch im Kapitel über Entlarvung schildert, zeigt, wie wissenschaftliche
Autorität als Verkaufskulisse dienen kann. Ethikkommissionen prüfen
Forschungsvorhaben mit Menschen, bevor sie beginnen. In der Medizin ist ihr
Votum vorgeschrieben; in der Psychologie und den Sozialwissenschaften
verlangen es viele Zeitschriften und Förderer. Für ein studentisches
Semesterprojekt genügt in der Regel die Abstimmung mit der Lehrperson;
sobald besonders schutzbedürftige Gruppen, sensible Daten, Täuschung oder
Belastungen im Spiel sind, ist die zuständige Stelle der Hochschule zu
fragen.}
\begin{description}
\item[{\rot\bf Interessenkonflikte}]~\\[-\lblskip]
\bi
  \ite Interessen hat jede und jeder -- entscheidend ist, sie offenzulegen
  \ite finanziell: Förderung, Beratung, Beteiligung
  \ite nicht-finanziell: eigene Theorie, persönliche Beziehungen,
    Überzeugungen
  \ite auch im Semesterprojekt: Befragen Sie Ihren eigenen Verein?
\ei
\os{\newpage}

\item[{\rot\bf Wann zur Ethikkommission?}]~\\[-\lblskip]
\bi
  \ite besonders schutzbedürftige Gruppen (Minderjährige, Patientinnen,
    Abhängige)
  \ite sensible Daten (Gesundheit, Religion, sexuelle Orientierung)
  \ite Täuschung, Belastung, Risiko für Teilnehmende
  \ite im Zweifel: früh fragen, nicht nachher
\ei
\os{\newpage}
\end{description}

%==============================================================================
\section{Grenzen der Wissenschaft}
%==============================================================================
\nsl{Nicht alles, was sich untersuchen lässt, darf untersucht werden, und
nicht jede Frage ist eine wissenschaftliche. Grenzen setzen zunächst die
Rechte der Beteiligten: Die Würde und Unversehrtheit der Teilnehmenden geht
dem Erkenntnisinteresse vor. Eine zweite Grenze betrifft Folgen außerhalb
der Studie, etwa wenn Ergebnisse missbraucht werden können
(\emph{Dual Use}). Eine dritte ist erkenntnistheoretisch: Wissenschaft kann
beschreiben, was ist, und vorhersagen, was unter Bedingungen geschieht. Ob
etwas geschehen \emph{soll}, folgt daraus nicht allein; dafür braucht es
Wertentscheidungen, die man offenlegen muss.}
\begin{description}
\item[{\rot\bf Was darf, was soll Wissenschaft nicht?}]~\\[-\lblskip]
\bi
  \ite \textbf{Schutz der Beteiligten}: Würde und Unversehrtheit gehen vor
  \ite \textbf{Folgen}: Ergebnisse, die missbraucht werden können
    (Dual Use)
  \ite \textbf{Sein und Sollen}: Daten zeigen, was ist -- nicht, was sein
    soll
  \ite Verantwortung endet nicht mit der Veröffentlichung
\ei
\os{\newpage}
\end{description}

\begin{merkbox}[Forschungsethik]
Der Zauberkünstler schützt seine Methode, die Forschung schützt ihre
Teilnehmenden. Offenlegen muss sie alles, was andere zur Prüfung brauchen
-- und geheim halten darf sie nur, was Menschen schützt.
\end{merkbox}
\os{\newpage}

%==============================================================================
\section{Übung in der Stunde: Ethik-Werkstatt}
%==============================================================================
\begin{description}
\item[{\rot\bf Ihr Semesterprojekt auf dem Prüfstand (25 Min.)}]~\\[-\lblskip]
\bi
  \ite Checkliste auf dem Übungsblatt Teil~A (Aufgabe~1): Was ist erfüllt, was
    offen, was tun Sie?
  \ite Einwilligungstext für Ihr Projekt anpassen (Aufgabe~2)
  \ite zu zweit: Fallvignetten sortieren -- FFP, fragwürdig oder in
    Ordnung? (Aufgabe~3)
\ei
\os{\newpage}
\end{description}

\begin{lehrhinweis}
Viele Semesterprojekte erheben keine sensiblen Daten; dann ist die
wichtigste Erkenntnis der Werkstatt, wie viel schon an Datenminimierung und
an der Trennung von Codeliste und Datensatz hängt. Projekte mit
Minderjährigen (etwa Schulklassen), Gesundheitsdaten oder Täuschung
gesondert besprechen und gegebenenfalls an die zuständige Stelle verweisen.
Der Erwartungshorizont zu den Fallvignetten steht am Ende von
Übungsblatt Teil~B.
\end{lehrhinweis}

%==============================================================================
\section{Verantwortung nach außen}
%==============================================================================
\nsl{Die Brücke zur Desinformationsabwehr schlägt das Buch mit einem
Vergleich. Der Zauberkünstler konstruiert ein Ereignis, das zu gut ist, um
wahr zu sein, und sagt offen, dass er täuscht. Stapel tat dasselbe und
verschwieg es. Der Unterschied liegt im Vertrag mit dem Publikum. Intransparenz
wird dort zur Täuschung, wo jemand den Eindruck erweckt, ein Verfahren sei
offen oder ein Befund geprüft, obwohl beides nicht stimmt. Das Buch warnt
zugleich vor der Nebenwirkung jedes Entlarvungstrainings: Wer gelernt hat,
dass alles inszeniert sein \emph{kann}, schließt leicht, dass alles
inszeniert \emph{ist} \parencite{mede2020populism}. Das Ziel ist
kalibriertes Vertrauen.}
\begin{description}
\item[{\rot\bf Wann wird Intransparenz zur Täuschung?}]~\\[-\lblskip]
\bi
  \ite \textbf{deklarierte} Täuschung: Zauberkunst -- alle wissen, dass
    getäuscht wird
  \ite \textbf{undeklarierte} Täuschung: Stapel -- Inszenierung, die sich
    als Prüfung ausgibt \parencite{levelt2012flawed}
  \ite Prüffragen: Wird mir Gewissheit verkauft, wo Glaubensgrade
    angemessen wären? Welche Alternative wird ausgeblendet?
  \ite Gegengift zum Zynismus: kalibriertes Vertrauen, nicht Generalverdacht
\ei
\os{\newpage}
\end{description}

%==============================================================================
\section{Rückblick: vom Staunen zur Verantwortung}
%==============================================================================
\nsl{Das Semester hat mit einem Pendel begonnen, das sich \enquote{von
allein} bewegte, und endet mit der Frage, was man über einen Trick verraten
darf. Dazwischen liegen dreizehn Staunmomente, aus denen jeweils ein
Begriff der Methodenlehre entstanden ist. Die Tabelle fasst sie zusammen.}
\begin{description}
\item[{\rot\bf Rückblick: von LE~1 (Staunen) bis LE~13 (Verantwortung)}]~\\[-\lblskip]
\nsl{\begin{center}\small
\begin{tabular}{rp{0.36\textwidth}p{0.44\textwidth}}
\toprule
\textbf{LE} & \textbf{Staunmoment} & \textbf{Kernbegriff} \\
\midrule
1 & Ideomotor-Pendel & Prediction Error \\
2 & Das geheime Objekt & Operationalisierung \\
3 & Kruskal Count & Suchstrategie \\
4 & Das Letzte Abendmahl & Apophenie/Bias \\
5 & Switchboard & Falsifikation \\
6 & Kontrollierte Seilroutine & Experimentelle Kontrolle \\
7 & Charpentier-Illusion & Confirmation Bias \\
8 & Chop Cup / Monty Hall & Bedingte Wahrscheinlichkeit \\
9 & Das Patter & Intersubjektivität \\
10 & Die unveränderliche Summe & Determinismus \\
11 & Patter-Umbau & Präzision \& Struktur \\
12 & CRED-Strategien & Credibility Cues \\
13 & Was verrät man? & Forschungsethik \\
\bottomrule
\end{tabular}
\end{center}
}
\os{\begin{center}
\begin{tabular}{rp{0.36\textwidth}p{0.44\textwidth}}
\toprule
\textbf{LE} & \textbf{Staunmoment} & \textbf{Kernbegriff} \\
\midrule
1 & Ideomotor-Pendel & Prediction Error \\
2 & Das geheime Objekt & Operationalisierung \\
3 & Kruskal Count & Suchstrategie \\
4 & Das Letzte Abendmahl & Apophenie/Bias \\
5 & Switchboard & Falsifikation \\
6 & Kontrollierte Seilroutine & Experimentelle Kontrolle \\
7 & Charpentier-Illusion & Confirmation Bias \\
\bottomrule
\end{tabular}
\end{center}
\newpage}
\os{\item[{\rot\bf Rückblick (Forts.)}]~\\[-\lblskip]
\begin{center}
\begin{tabular}{rp{0.36\textwidth}p{0.44\textwidth}}
\toprule
\textbf{LE} & \textbf{Staunmoment} & \textbf{Kernbegriff} \\
\midrule
8 & Chop Cup / Monty Hall & Bedingte Wahrscheinlichkeit \\
9 & Das Patter & Intersubjektivität \\
10 & Die unveränderliche Summe & Determinismus \\
11 & Patter-Umbau & Präzision \& Struktur \\
12 & CRED-Strategien & Credibility Cues \\
13 & Was verrät man? & Forschungsethik \\
\bottomrule
\end{tabular}
\end{center}
}
\os{\newpage}

\item[{\rot\bf Der Bogen des Semesters}]~\\[-\lblskip]
\begin{center}
\resizebox{0.95\textwidth}{!}{%
\begin{tikzpicture}[font=\sffamily,>=Stealth,
  st/.style={draw=rwuviolet,thick,fill=rwuviolet!10,rounded corners=3pt,
    minimum height=14mm,align=center,text width=30mm}]
  \node[st] (a) at (0,0) {\textbf{Staunen}\\LE 1--2};
  \node[st] (b) at (4,0) {\textbf{Suchen, prüfen}\\LE 3--7};
  \node[st] (c) at (8,0) {\textbf{Messen, deuten}\\LE 8--10};
  \node[st] (d) at (12,0) {\textbf{Mitteilen}\\LE 11--12};
  \node[st,draw=rwucyan,fill=rwucyan40!40] (e) at (16,0) {\textbf{Verantworten}\\LE 13};
  \draw[->,very thick] (a) -- (b);
  \draw[->,very thick] (b) -- (c);
  \draw[->,very thick] (c) -- (d);
  \draw[->,very thick] (d) -- (e);
  \draw[->,thick,rwugray,dashed] (e.south) to[bend left=12] node[below]{neue Befunde erzeugen neues Staunen} (a.south);
\end{tikzpicture}}
\end{center}
\bi
  \ite Staunen erzeugen -- Erklärung suchen -- Methode verstehen, und
    die Methode verantworten
\ei
\os{\newpage}
\end{description}

%==============================================================================
\section{Lernziele und Aufgaben}
%==============================================================================
\begin{description}
\item[{\rot\bf Nach dieser Einheit können Sie}]~\\[-\lblskip]
\bi
  \ite ethische Anforderungen an Ihr eigenes Forschungsvorhaben benennen
  \ite eine informierte Einwilligung formulieren
  \ite Anonymisierung und Pseudonymisierung unterscheiden
  \ite Fabrication, Falsification und Plagiarism von fragwürdigen und
    redlichen Praktiken abgrenzen
  \ite begründen, was in der Forschung offengelegt werden muss und was
    geschützt werden darf
\ei
\os{\newpage}

\item[{\rot\bf Zum Abschluss}]~\\[-\lblskip]
\bi
  \ite \textbf{Ethik-Checkliste} für Ihr Semesterprojekt fertigstellen
    (Aufgabe~1); offene Punkte vor der Datenerhebung klären
  \ite \textbf{Abschluss-Tagebucheintrag} (Aufgabe~5): Was hat Sie in
    diesem Semester am meisten überrascht -- und was hat das mit
    Wissenschaft zu tun?
  \ite Tagebuch des Staunens zur Portfolioprüfung abgeben
  \ite Seminararbeit: Ethik und Datenumgang im Methodenteil beschreiben
\ei
\os{\newpage}
\end{description}

\handoutinput{handouts/LE13}
